
This work builds on and extends three lines of research: benchmark generalization studies, underspecification in machine learning, and texture bias analysis.

\subsubsection{Benchmark Generalization Studies}

Recht et al. (2019) provided the first systematic evidence that ImageNet classifiers do not generalize to ImageNet-like test sets. By creating ImageNetV2 with a carefully matched distribution, they demonstrated consistent 10-15\% accuracy drops across 70+ models. Subsequent work extended this analysis to CIFAR-10 with the CINIC-10 dataset \citep{darlow2018cinic}, finding similar generalization gaps. However, these studies examine individual datasets in isolation without connecting generalization failure to dataset popularity or the broader ML ecosystem's optimization patterns.

Engstrom et al. (2020) and Taori et al. (2020) further documented ''effective robustness'' differences across model families, showing that ImageNet accuracy improvements do not uniformly translate to out-of-distribution performance. The present work extends this line by proposing that \textit{popularity itself} is a predictor of generalization failure.

\subsubsection{Underspecification and Benchmark Overfitting}

D'Amour et al. (2020) introduced the underspecification framework, demonstrating that models achieving identical benchmark performance can exhibit dramatically different behavior under distribution shift. Their theoretical analysis explains why benchmark performance fails to predict deployment success: the optimization surface admits many equivalent solutions, and benchmarks do not constrain which one the model learns. This provides theoretical grounding for the hypothesis tested here: if popular benchmarks receive more optimization pressure, models may converge to increasingly benchmark-specific solutions.

The benchmark co-evolution concept also relates to Goodhart's Law in ML \citep{thomas2022reliance}---when a measure becomes a target, it ceases to be a good measure. This work operationalizes this concern by measuring the actual research investment differential between high-use and low-use datasets.

\subsubsection{Texture Bias and Feature Learning}

Geirhos et al. (2019) demonstrated that ImageNet-trained CNNs rely heavily on texture rather than shape, contrary to human perception. Hermann et al. (2020) traced this bias to training data statistics, suggesting that benchmark-specific artifacts shape learned representations. The hypothesis that intensive optimization on popular benchmarks might amplify texture bias was tested as a mechanism for the co-evolution effect.

However, the experiments contradict this expectation at CIFAR scale: ResNet-18 shows \textit{lower} texture bias (0.126) than VGG-11 (0.161). This suggests that architectural innovations (skip connections, batch normalization) may counteract texture exploitation, and that ImageNet-scale findings do not directly transfer to 32x32 image classification.

